\documentclass[12pt, a4paper]{article}
\usepackage[utf8]{inputenc}
\usepackage[T1]{fontenc}
\usepackage[french]{babel}
\usepackage[margin=2.5cm]{geometry}
\usepackage{amsmath, amssymb, amsthm, mathtools}
\usepackage{listings}
\usepackage{xcolor}
\usepackage{tcolorbox}
\tcbuselibrary{theorems, skins, breakable}
\usepackage{booktabs}
\usepackage{fancyhdr}
\usepackage{hyperref}
\usepackage{tikz}
\usetikzlibrary{arrows.meta, positioning}

\definecolor{suva_blue}{RGB}{30, 100, 180}
\definecolor{code_bg}{RGB}{245, 247, 250}
\definecolor{code_kw}{RGB}{0, 80, 200}
\definecolor{code_str}{RGB}{180, 40, 40}
\definecolor{code_cmt}{RGB}{80, 130, 80}
\definecolor{warn_bg}{RGB}{255, 248, 220}
\definecolor{success_bg}{RGB}{230, 255, 235}

\lstdefinestyle{python}{
  language=Python,
  basicstyle=\ttfamily\small,
  keywordstyle=\color{code_kw}\bfseries,
  stringstyle=\color{code_str},
  commentstyle=\color{code_cmt}\itshape,
  backgroundcolor=\color{code_bg},
  frame=single, rulecolor=\color{gray!40},
  numbers=left, numberstyle=\tiny\color{gray},
  breaklines=true, tabsize=4, showstringspaces=false,
}

\newtcolorbox{defbox}[2][]{colback=success_bg,colframe=green!50!black,
  fonttitle=\bfseries,title={Définition : #2},#1}
\newtcolorbox{theorembox}[2][]{colback=blue!5,colframe=suva_blue,
  fonttitle=\bfseries,title={Théorème : #2},#1}
\newtcolorbox{exercice}[2][]{colback=white,colframe=suva_blue!60,
  fonttitle=\bfseries,title={Exercice #2},#1,breakable}
\newtcolorbox{solution}[1][]{colback=gray!5,colframe=gray!40,
  fonttitle=\bfseries\itshape,title={Solution},#1,breakable}
\newtcolorbox{importantbox}[1][]{colback=suva_blue!10,colframe=suva_blue,
  fonttitle=\bfseries,title={Connexion avec le projet},#1}
\newtcolorbox{warningbox}[1][]{colback=warn_bg,colframe=orange!80!black,
  fonttitle=\bfseries,title={Attention},#1}

\pagestyle{fancy}
\fancyhf{}
\fancyhead[L]{\color{suva_blue}\textbf{SuVa --- Chapitre 4}}
\fancyhead[R]{\color{gray}Réseaux de Treillis}
\fancyfoot[L]{\footnotesize\color{gray}Mohamed Lamine MBENGUE}
\fancyfoot[C]{\thepage}
\fancyfoot[R]{\footnotesize\color{gray}portfolio-alamine.web.app}
\renewcommand{\headrulewidth}{0.4pt}
\renewcommand{\footrulewidth}{0.3pt}

% ══════════════════════════════════════════════════════════════════════
\begin{document}

\begin{center}
  {\color{suva_blue}\rule{\linewidth}{2pt}}\\[0.5cm]
  {\Huge\bfseries\color{suva_blue} Chapitre 4}\\[0.3cm]
  {\large\bfseries Réseaux de Treillis et Cryptographie Post-Quantique}\\[0.2cm]
  {\small\color{gray} La sécurité de Dilithium repose sur des problèmes difficiles sur les réseaux}\\[0.3cm]
  {\color{suva_blue}\rule{\linewidth}{2pt}}
\end{center}

\vspace{0.3cm}
\begin{tcolorbox}[colback=suva_blue!5, colframe=suva_blue!40,
  left=6pt, right=6pt, top=4pt, bottom=4pt]
\small
\begin{tabular}{ll}
  \textbf{Auteur}      & Mohamed Lamine MBENGUE \\
  \textbf{Spécialités} & Sécurité · DevOps · Dev Full Stack Web \& Mobile \\
  \textbf{Contact}     & +221 78 178 44 36 · mbenguemohamedlamine2022@gmail.com \\
  \textbf{Portfolio}   & \href{https://portfolio-alamine.web.app}{\texttt{portfolio-alamine.web.app}} \\
\end{tabular}
\end{tcolorbox}
\vspace{0.3cm}

\tableofcontents
\newpage

%══════════════════════════════════════════════════════════════════════
\section{Pourquoi la crypto actuelle sera cassée par les ordinateurs quantiques}
%══════════════════════════════════════════════════════════════════════

\subsection{ECDSA et la vulnérabilité quantique}

Ethereum utilise \textbf{ECDSA} (Elliptic Curve Digital Signature Algorithm).
Sa sécurité repose sur la difficulté du \textbf{problème du logarithme
discret sur les courbes elliptiques}~:

\begin{center}
Connaissant $P$ (point) et $Q = k \cdot P$, trouver $k$.
\end{center}

\textbf{L'algorithme de Shor} (1994) résout ce problème en temps polynomial
sur un ordinateur quantique. Conséquence~:
\begin{itemize}
  \item Un QC suffisant peut retrouver la clé privée depuis la clé publique
  \item Toutes les transactions Ethereum deviendraient forgeable
  \item $\sim$\$200 milliards de stablecoins exposés
\end{itemize}

\subsection{Pourquoi les réseaux résistent au quantique}

Les meilleurs algorithmes quantiques connus (Grover, Shor) n'offrent pas
d'avantage polynomial pour les problèmes sur les réseaux.

L'algorithme de Shor \textbf{ne fonctionne pas} sur les problèmes LWE/SIS.

%══════════════════════════════════════════════════════════════════════
\section{Qu'est-ce qu'un réseau de treillis (lattice) ?}
%══════════════════════════════════════════════════════════════════════

\begin{defbox}{Réseau de treillis}
Un réseau $\Lambda$ de rang $n$ est l'ensemble des combinaisons
\textbf{entières} de $n$ vecteurs linéairement indépendants $b_1, \ldots, b_n$~:
$$\Lambda = \left\{ \sum_{i=1}^n a_i b_i \mid a_i \in \mathbb{Z} \right\}$$
La matrice $B = [b_1 | \cdots | b_n]$ est appelée \textbf{base} du réseau.
\end{defbox}

\textbf{Visualisation 2D~:}

\begin{center}
\begin{tikzpicture}[scale=0.8]
  % Points du reseau
  \foreach \i in {-2,-1,0,1,2} {
    \foreach \j in {-2,-1,0,1,2} {
      \fill (\i*2 + \j*0.5, \j*1.5) circle (2pt);
    }
  }
  % Base vectors
  \draw[->, thick, suva_blue] (0,0) -- (2,0) node[right] {$b_1$};
  \draw[->, thick, red] (0,0) -- (0.5,1.5) node[above] {$b_2$};
  % Axes
  \draw[gray, ->] (-4,0) -- (5.5,0);
  \draw[gray, ->] (0,-3) -- (0,4);
\end{tikzpicture}
\end{center}

%══════════════════════════════════════════════════════════════════════
\section{Problèmes difficiles sur les réseaux}
%══════════════════════════════════════════════════════════════════════

\subsection{SVP --- Shortest Vector Problem}

\begin{defbox}{SVP}
\textbf{Problème du vecteur le plus court~:}
Trouver le vecteur non-nul le plus court dans un réseau $\Lambda$.
$$\lambda_1(\Lambda) = \min_{v \in \Lambda \setminus \{0\}} \|v\|$$
\end{defbox}

SVP est NP-difficile dans le cas général. Les meilleurs algorithmes
(BKZ, LLL) ont une complexité exponentielle.

\subsection{LWE --- Learning With Errors}

\begin{defbox}{LWE}
\textbf{Problème LWE~:} Connaissant les couples $(a_i, b_i)$ où~:
$$b_i = \langle a_i, s \rangle + e_i \pmod{q}$$
avec $s$ le \textbf{secret}, $a_i$ aléatoires uniformes, et $e_i$ un
\textbf{bruit} petit, retrouver $s$.
\end{defbox}

\begin{importantbox}
\textbf{Analogie~:} Imagine qu'on te donne 100 équations du type~:
$$a_1 x_1 + a_2 x_2 + \cdots + a_n x_n + e \equiv b \pmod{q}$$
Chaque équation a une petite erreur $e$. Sans les erreurs, le système
est facilement résoluble (Gauss). \textbf{Avec les erreurs}, même avec
des millions d'équations, retrouver $(x_1, \ldots, x_n)$ est
computationnellement infaisable.
\end{importantbox}

\subsection{MLWE --- Module Learning With Errors}

\begin{defbox}{MLWE}
Version modulaire de LWE~: le secret $s$ et le bruit $e$ sont des
\textbf{vecteurs de polynômes} dans $R_q^l$, et $A$ est une matrice
$k \times l$ de polynômes dans $R_q$.

$$t = A s_1 + s_2 \pmod{q}$$

où $s_1 \in R_q^l$ et $s_2 \in R_q^k$ ont des coefficients petits
(bornés par $\eta$).
\end{defbox}

\begin{importantbox}
Dans Dilithium~:
\begin{itemize}
  \item $A \in R_q^{4 \times 4}$ est publique (générée de $\rho$)
  \item $s_1 \in R_q^4$ a des coefficients dans $[-\eta, \eta] = [-2, 2]$ (secret)
  \item $s_2 \in R_q^4$ a des coefficients dans $[-\eta, \eta] = [-2, 2]$ (secret)
  \item $t = As_1 + s_2$ est publique (dans la clé publique)
\end{itemize}
\textbf{Sécurité~:} Étant donné $(A, t)$, retrouver $(s_1, s_2)$ est MLWE.
\end{importantbox}

\subsection{MSIS --- Module Short Integer Solution}

\begin{defbox}{MSIS}
\textbf{Module Short Integer Solution~:}
Étant donné $A \in R_q^{k \times l}$, trouver un vecteur court non-nul
$z \in R_q^l$ tel que~:
$$Az \equiv 0 \pmod{q}$$
\end{defbox}

\begin{importantbox}
\textbf{Lien avec Dilithium~:} L'unforgeability (impossibilité de forger
une signature) repose sur MSIS. Un attaquant qui forge une signature
produit essentiellement un court vecteur $z$ tel que $Az \approx ct_1 \cdot 2^d$.
Cela serait une solution au problème MSIS.
\end{importantbox}

%══════════════════════════════════════════════════════════════════════
\section{L'attaque BKZ --- comment mesurer la sécurité}
%══════════════════════════════════════════════════════════════════════

\begin{defbox}{Algorithme BKZ}
BKZ (Block Korkine-Zolotarev) est le meilleur algorithme pratique
pour réduire une base de réseau (trouver de courts vecteurs).
Il paramétrisé par un \textbf{block size} $\beta$~:

\begin{itemize}
  \item Plus $\beta$ est grand, plus le résultat est bon
  \item Mais le temps d'exécution est $\sim 2^{0.292\beta}$ (classique)
  \item Et $\sim 2^{0.265\beta}$ (quantique, algorithme de Grover)
\end{itemize}

La sécurité en bits est $\approx 0.292\beta$.
\end{defbox}

\subsection{Estimation de sécurité pour Dilithium2}

Les scripts dans \texttt{security-estimates-simon/} calculent le
$\beta$ nécessaire pour casser le schéma~:

\begin{lstlisting}[style=python]
# Simplifie depuis Falcon_MLWE_security.py
def estimate_mlwe_security(n, q, k, l, eta):
    """
    Estime la securite de l'instance MLWE.
    Trouve le block-size BKZ minimal pour casser le schema.
    """
    # Dimension du reseau : (k+l)*n
    lattice_dim = (k + l) * n

    # Cherche le beta optimal
    for beta in range(50, 1000):
        # Longueur du vecteur cible (approximation)
        # Pour trouver s1, s2 avec coefficients <= eta
        target_norm = eta * (lattice_dim ** 0.5)

        # Longueur minimale que BKZ peut trouver (Hermite factor)
        hermite_factor = (
            (beta / (2 * 3.14159 * 2.71828)) **
            (1 / (2 * beta)) *
            (q ** (k * n / lattice_dim))
        )
        expected_norm = hermite_factor ** lattice_dim

        if expected_norm <= target_norm:
            # Ce beta suffit pour l'attaque
            security_classical = 0.292 * beta
            security_quantum = 0.265 * beta
            return beta, security_classical, security_quantum

    return None

# Dilithium2
q = 8_380_417
n = 256; k = 4; l = 4; eta = 2
beta, sec_c, sec_q = estimate_mlwe_security(n, q, k, l, eta)
print(f"Block size BKZ : {beta}")
print(f"Securite classique : {sec_c:.0f} bits")
print(f"Securite quantique : {sec_q:.0f} bits")
# Attendu : ~426 bits, ~124 bits classique, ~112 bits quantique
\end{lstlisting}

%══════════════════════════════════════════════════════════════════════
\section{SelfTargetMSIS --- sécurité de la signature}
%══════════════════════════════════════════════════════════════════════

La sécurité de Dilithium repose sur un problème encore plus fort~:
le \textbf{SelfTargetMSIS}. Forger une signature sans la clé secrète
revient à trouver un vecteur court $(z, c)$ tel que~:
$$Az - ct_1 \cdot 2^d \approx 0 \pmod{q}$$
avec $\|z\|_\infty < \gamma_1 - \beta$ et $c$ un polynôme sparse.

C'est encore plus difficile que MSIS car le target $ct_1 \cdot 2^d$ dépend
du message à signer.

%══════════════════════════════════════════════════════════════════════
\section{Exercices}
%══════════════════════════════════════════════════════════════════════

\begin{exercice}{1 --- LWE jouet}
Soit $q = 101$, $n = 3$, secret $s = (2, 7, 1)$, bruit max $|e| \leq 1$.

On te donne ces échantillons LWE~:
\begin{align*}
  (a_1, b_1) &= ((3, 1, 4), 3 \cdot 2 + 1 \cdot 7 + 4 \cdot 1 + e_1)\\
  (a_2, b_2) &= ((1, 5, 9), 1 \cdot 2 + 5 \cdot 7 + 9 \cdot 1 + e_2)
\end{align*}

\begin{enumerate}[label=\alph*)]
  \item Calcule les $b_i$ exacts (sans erreur) en arithmétique $\mathbb{Z}_{101}$.
  \item Ajoute un bruit $e_1 = 1$, $e_2 = -1$ et calcule les $b_i$ bruitées.
  \item Si tu as accès à plein d'échantillons propres (sans bruit), peux-tu
        récupérer $s$ facilement ? (Indice~: Gauss)
  \item Avec le bruit, pourquoi est-ce difficile ?
\end{enumerate}
\end{exercice}

\begin{solution}
\begin{lstlisting}[style=python]
q = 101
s = [2, 7, 1]

a1, a2 = [3, 1, 4], [1, 5, 9]

# a) Sans bruit
b1_exact = sum(a*si for a, si in zip(a1, s)) % q
b2_exact = sum(a*si for a, si in zip(a2, s)) % q
print(f"b1 exact = {b1_exact}")   # (6+7+4) = 17
print(f"b2 exact = {b2_exact}")   # (2+35+9) = 46

# b) Avec bruit
b1_noisy = (b1_exact + 1) % q
b2_noisy = (b2_exact - 1) % q
print(f"b1 bruite = {b1_noisy}")   # 18
print(f"b2 bruite = {b2_noisy}")   # 45

# c) Sans bruit : on peut faire Gauss.
# Avec seulement 3 variables et 3 equations propres :
# 3x + y + 4z = 17
# x + 5y + 9z = 46
# ... -> solution unique = (2, 7, 1)

# d) Avec bruit : b = As + e.
# On ne sait pas quels echantillons ont quel bruit.
# Gauss donne une solution fausse.
# Les meilleurs algorithmes prennent un temps exponentiel.
\end{lstlisting}
\end{solution}

\begin{exercice}{2 --- MLWE dans Dilithium}
Dans Dilithium2 ($k=l=4$, $n=256$, $q=8\,380\,417$, $\eta=2$)~:
\begin{enumerate}[label=\alph*)]
  \item Quelle est la dimension totale du réseau utilisé pour l'attaque ?
  \item Si on augmentait $\eta$ de 2 à 4, que se passe-t-il pour la sécurité ?
  \item Si on augmentait $q$ (corps plus grand), que se passe-t-il ?
\end{enumerate}
\end{exercice}

\begin{solution}
\begin{lstlisting}[style=python]
# a) Dimension du reseau
n, k, l = 256, 4, 4
dim = (k + l) * n
print(f"Dimension : {dim}")  # 2048

# b) eta plus grand => vecteur secret plus long => plus facile a trouver
# => securite DIMINUE si on augmente eta
# C'est pourquoi Dilithium3 augmente k, l plutot qu'eta.

# c) q plus grand => le reseau est "plus facile" a reduire proportionnellement
# => securite DIMINUE si on augmente q avec les memes k, l.
# C'est le probleme du ZK-Dilithium avec BabyBear (q=2^31-1) :
# plus grand q => il faut augmenter k, l pour compenser.
\end{lstlisting}
\end{solution}

\begin{exercice}{3 --- Explorer les scripts de sécurité}
Ouvre \texttt{security-estimates-simon/Falcon\_Dilithium.py}.
\begin{enumerate}[label=\alph*)]
  \item Quelle est la sécurité estimée pour Dilithium2 en bits (classique et quantique) ?
  \item Comment la sécurité est-elle estimée (quelle méthode, quel modèle) ?
  \item Quelle différence entre attaque MLWE et MSIS dans le contexte de Dilithium ?
\end{enumerate}
\end{exercice}

\end{document}
